\section{Methods}
\label{sec:methods}

\subsection{RL Problem Definition}
\label{RL_Def}

\begin{figure*}[t]
   \centering
   \includegraphics[width=\linewidth]{figures/arch.png}
   \caption{\textbf{Hierarchical structure of Hestia.} Hestia first predicts the camera's look-at point $L_t$ using a proposal neural network that takes grid information $G_t$ processed from the depth image $D_t$ and the camera pose as input. Next, Hestia employs a grid encoder to encode the grid information $G_t$ and performs trilinear interpolation to extract corresponding features from the encoded grid at different layers based on $L_t$. These multilevel interpolated features are then concatenated with the vector information $M_t$ which includes the camera pose $X_t$ and the maximum flyable height, $H_t$ as well as the encoded image features. The image features are extracted using an image encoder, which takes the grayscale image $I_t$ as input. Finally, this combined feature representation is fed into the RL policy model to predict the camera's position $a_t$. Note that Hestia adopts $a'_t$, the nearest collision-free point to $a_t$, as the final camera position to ensure a collision-free viewpoint. Hence, the next-best viewpoint $\{a'_t, L_t\}$ is used for data collection.}
   \label{fig:arch}
\end{figure*}

Our next-best-view task is to identify a 5-DoF viewpoint that maximizes the coverage ratio of an incomplete occupancy grid of the scene. 
The task's goal is similar to greedy methods, which always seek the locally optimal solution.
We formulate the problem as a Markov Decision Process (MDP), denoted by the tuple $\{S, A, P, R, \gamma\}$. 
At each time step $t$, the agent with an RGB-D camera observes a state $s_t$ from the set of all possible states $S$ and chooses an action $a_t$ from the action space $A$. 
The environment then transitions to the subsequent state $s_{t+1}$ according to the probabilities described by $P$, and provides a reward $r_t$.
The magnitude of this reward is determined by the reward function: %$R(\cdot \mid s, a): S \times A \to r$.
\begin{equation}\label{eq:reward_function}
     R(\cdot \mid s, a): S \times A \to r.
\end{equation}
In reinforcement learning, the main goal is to discover an optimal policy $\pi$ that maximizes the expected sum of discounted rewards, given by: %$E_t = \sum_{k=0}^{\infty} \gamma^k r_{t+k},$
\begin{equation}\label{eq:discounted_reward}
     E_t = \sum_{k=0}^{\infty} \gamma^k r_{t+k},
\end{equation}
where $\gamma \in (0, 1)$ is the discount factor.
We set $\gamma$ to 0.1 to align with the greedy-like objective and to avoid spurious correlations from large positive future rewards (see Fig.~S8).
\\
\textbf{State space.} The state space of Hestia is defined as: %$S = \Big\{ s_t \;\Big|\; s_t = \big\{ I_t, M_t, G_t, L_t \big\}, \; t \in \mathbb{N} \Big\}, $
\begin{equation}\label{eq:state_space_hestia}
 S = \Big\{ s_t \;\Big|\; s_t = \big\{ I_t, M_t, G_t, L_t \big\}, \; t \in \mathbb{N} \Big\}
 \end{equation}
where $I_t \in \mathbb{R}^{h \times w}$ is the grayscale image with height $h$ and width $w$, and $L_t \in \mathbb{R}^{3}$ is the camera look-at point. 
The vector $M_t \in \mathbb{R}^{6}$ consists of $X_t \in \mathbb{R}^{5}$, which is the camera position, pitch, and yaw, as well as $H_t \in \mathbb{R}^{1}$, representing the maximum flyable height for the capture. 
Meanwhile, $G_t \in \mathbb{R}^{g \times g \times g \times 10}$ includes the aggregated grid information at resolution $g$, consisting of $O_t \in \mathbb{R}^{g \times g \times g \times 1}$ for the cumulative occupancy grid, $C_t \in \mathbb{R}^{g \times g \times g \times 3}$ for the positional encoding, and $F_t \in \{0, 1\}^{g \times g \times g \times 6}$ for the cumulative face visibility.
The cumulative face visibility is updated iteratively as: %$F_t = f_t \lor F_{t-1},$
\begin{equation}
 F_t = f_t \lor F_{t-1}
\end{equation}  
where $F_t$ represents the cumulative face visibility for all voxels up to time $t$, and $f_t \in \{0, 1\}^{g \times g \times g \times 6}$ denotes the current face visibility.
%denotes the face visibility for the current timestep $t$. 
To compute $f_t$, the depth image $D_t$ is unprojected into a voxelized point cloud $V = \{v_i \;|\; i \in \mathbb{N}\}$, where $v_i$ is the $i$-th voxel. 
Each voxel $v_i$ is associated with a viewing direction vector $d_{v_i} \in \mathbb{R}^3$, defined as the vector pointing from the voxel center to the collision-free camera position $a^{'}_{t}$. 
The vector $\mathbf{d}_{v_i}$ is computed as: %$ \mathbf{d}_{v_i} = \frac{a^{'}_{t} - p_{v_i}}{\|a^{'}_{t} - p_{v_i}\|},$
\begin{equation}
\mathbf{d}_{v_i} = \frac{a^{'}_{t} - p_{v_i}}{\|a^{'}_{t} - p_{v_i}\|}
\end{equation}  
where $p_{v_i}$ is the center of voxel $v_i$.
For each voxel $v_i$ and its six outward-facing face normals $n_{i,j} \in \mathbb{R}^3$, the face visibility is determined and aggregated as
\begin{align}
f_t(v_i, j) &= \mathbbm{1}\!\left(d_{v_i} \cdot n_{i,j} > 0\right), \notag \\
&\quad \forall v_i \in V, \; j \in \{1, \dots, 6\}
\end{align}
where $\mathbbm{1}(\cdot)$ is the indicator function.
By iterating over all voxels and their respective faces, $f_t$ is constructed, and the cumulative visibility $F_t$ is updated accordingly.
Although this method cannot handle all face visibilities, the approximation enables efficient computation of face visibility.
Moreover, non-visible voxel faces simply contribute no reward and therefore do not affect the next-best-view selection.
Unlike prior works~\cite{chen2024gennbv}, which consider only $O_t$ and $C_t$ and thereby treat voxels as points, we treat voxels as cubes to mitigate the information loss caused by approximating voxels as points (see~\cref{fig:voxel_ray}).
For details regarding $O_t$ and $C_t$, please refer to the work~\cite{chen2024gennbv}.
\\
\textbf{Action space.} The action space: %$A = \Bigl\{ a_t \;\Big|\; a_t \in [-1, 1]^3, \; t \in \mathbb{N} \Bigr\}$
\begin{equation}
A = \Bigl\{ a_t \;\Big|\; a_t \in [-1, 1]^3, \; t \in \mathbb{N} \Bigr\}
\end{equation}
represents the set of possible 3-DoF viewpoints (e.g., camera positions) at each time step $t$, where each coordinate is initially bounded within $[-1, 1]$. 
These coordinates are subsequently normalized to the environment’s scale to ensure appropriate positioning within the scene.
Additionally, the camera’s pitch and yaw are derived from the look-at point and the collision-free action $a^{'}_{t}$ converted from $a_t$ (see~\cref{sec:nbv_net}). %\\ \\
\\
\textbf{Reward. }
The reward function is defined as:  %$r_t = R(s_t, a_t) = r_{\text{coverage}}(s_t, a_t) + r_{\text{constraint}}(s_t, a_t),$
\begin{equation}
r_t = R(s_t, a_t) = r_{\text{coverage}}(s_t, a_t) + r_{\text{constraint}}(s_t, a_t)
\end{equation}
where $r_{\text{coverage}}(s_t, a_t)$ encourages the observation of new voxel faces and is expressed as:  
\begin{align}
r_{\text{coverage}}(s_t, a_t) 
&= \frac{\sum_{i=1}^N \sum_{j=1}^6 \left( F_t^{i,j} - F_{t-1}^{i,j} \right) \cdot M_{\text{col}}}{N \cdot 6} \notag \\
&\quad \cdot 0.3
\end{align}
where $F_t^{i,j}$ and $F_{t-1}^{i,j}$ represent the visibility status of the $j$-th face of the $i$-th voxel at time $t$ and $t-1$, respectively. Here, $M_{\text{col}} \in \{0, 1\}$ is a collision indicator, set to $0$ in the event of a collision, thereby preventing any positive reward for invalid actions.
%The term $r_{\text{constraint}}(s_t, a_t)$ penalizes unsafe and invalid actions and is defined as
The term $r_{\text{constraint}}(s_t, a_t) = -0.01$ is applied when unsafe or invalid actions occur (see~\cref{asec:rew} for details).
% \begin{equation}
% r_{\text{constraint}}(s_t, a_t) =
% \begin{cases} 
% -0.01, & \text{if } r_{\text{coverage}}(s_t, a_t) = 0, \\
%          & \text{or } a_t[2] > H_t, \\
%          & \text{or } a_t \in \text{non-free voxels}, \\
% 0,       & \text{otherwise.}
% \end{cases}
% \end{equation}
Our reward is based on the face coverage ratio rather than the point coverage ratio to ensure more comprehensive capture (see~\cref{fig:voxel_ray}).  
Furthermore, to prevent spurious correlations, the reward design aligns with a greedy-like objective, which differs significantly from prior works~\cite{chen2024gennbv, peralta2020next} that provide a large goal reward when the coverage ratio reaches a predefined target.

\subsection{Next-Best-View Hierarchical Network}
\label{sec:nbv_net}

The goal of the task is to predict a 5-DoF viewpoint for data collection.  
Directly modeling the 5-DoF viewpoint in the RL continuous action space is challenging due to the high-dimensional search space.  
To address this, Hestia introduces a hierarchical structure to simplify the problem. %\\ \\  
\\
\textbf{Look-at point prediction.} Hestia first predicts the look-at point using a proposal network (see~\cref{fig:arch}), which takes grid information $G_t$ as input to determine \textit{where to look}.  
The proposal network is a 3D convolutional neural network with a self-attention layer to expand the receptive field.
% \begin{figure*}[t]
%   \centering
%   \includegraphics[width=\linewidth]{figures/system.png}
%   \caption{\textbf{The proposed data collection system and flowchart.} Hestia goes beyond prior works by introducing a prototype of an intelligent data collection system utilizing a drone equipped with an RGB camera for data capture, HTC Lighthouse base stations and Crazyflie 2.1 for localization, and a Wifi router for wireless communication.}
%   %Although tested indoors, this system highlights the potential of drones with next-best-view planners for large-scale data collection.}
%   \label{fig:real_world}
% \end{figure*}
The output is then passed through linear layers to decode the look-at point $L_t$.  
To model the look-at point as a probability distribution, the reparameterization trick is used, treating it as a sample from a normal distribution.  %\\ \\
\\
\textbf{Viewpoint position prediction. } To predict the remaining 3-DoF viewpoint position (e.g., \textit{where to fly}), the grid information is encoded into a multilevel feature grid using a shallow 3D CNN.  
The look-at point $L_t$ is then used to perform trilinear interpolation on the multilevel features from the grid.  
These interpolated features are concatenated with the image embedding, which is extracted by an image encoder, a shallow CNN that takes the grayscale image $I_t$ as input.  
Additionally, the features are concatenated with vector information $M_t$, which includes the camera pose $X_t$ and the maximum flyable height $H_t$.
The combined features are fed into the RL policy model to predict the action $a_t$.  
While the reward function helps constrain $a_t$ to avoid collisions, an additional constraint is applied to ensure a collision-free viewpoint.  
Specifically, $a_t$ is shifted to its nearest collision-free point $a^{'}_{t}$ determined using $G_t$ and $H_t$.  
This adjusted action $a^{'}_{t}$ serves as the final viewpoint position for data capture.  
Thus, Hestia's next-best-view is represented as $\{a^{'}_{t}, L_{t}\}$. %\\ \\
\\
\textbf{Training loss functions.} %Fully joint learning of two continuous action spaces through RL is complex and unnecessary for our task. 
The look-at point and viewpoint prediction networks can be jointly trained using the \textit{RL reward} due to their connection. 
However, our previous experiments showed no clear benefit from joint training using the \textit{RL reward}. 
To simplify the design, we detach the gradient flow between the networks and train the look-at point prediction network with supervised learning using ground-truth targets.
The entire architecture of Hestia in~\cref{fig:arch} is trained together without any pretraining on other datasets or tasks.
The ground truth look-at point $L_t^\text{gt}$ is computed as the weighted average position of the ground truth uncaptured surface: %$L_t^{\text{gt}} = \frac{\sum_{v_i \in U} w_{v_i} \, p_{v_i}}{\sum_{v_i \in U} w_{v_i}}, $
\begin{equation}
L_t^{\text{gt}} = \frac{\sum_{v_i \in U} w_{v_i} \, p_{v_i}}{\sum_{v_i \in U} w_{v_i}}
\end{equation}
where $U$ represents the set of voxels containing ground truth uncaptured faces, and $w_{v_i}$
%$w_{v_i} = \sum_{f \in F_{v_i}^{\text{gt}}} 1$ 
is defined as the total number of ground truth uncaptured faces within voxel $v_i$:
\begin{equation}
w_{v_i} = \sum_{f \in F_{v_i}^{\text{gt}}} 1
\end{equation}
where $F_{v_i}^{\text{gt}}$ is the set of ground truth uncaptured faces associated with voxel $v_i$.  
Thus, the loss function for the proposal network is formulated as: %$\mathcal{L}_{\text{proposal}} = \|L_t -L_t^{\text{gt}}\|^2.$
\begin{equation}
\mathcal{L}_{\text{proposal}} = \|L_t -L_t^{\text{gt}}\|^2
\end{equation}
The loss of the viewpoint prediction network is the same as the regular RL loss $\mathcal{L}_{\text{RL}}$ which depends on the RL method used (see~\cref{asec:train_det}), combined with an auxiliary loss: %$\mathcal{L}_{\text{aux}} = \|a_t - a^{'}_t\|^2$
\begin{equation}
\mathcal{L}_{\text{aux}} = \|a_t - a^{'}_t\|^2
\end{equation}
to encourage the predicted action $a_t$ to align with the collision-free action $a^{'}_t$.  
Hence, the overall loss function for Hestia is %$\mathcal{L}_\text{all} = \mathcal{L}_{\text{RL}} + 0.5 \cdot \mathcal{L}_{\text{aux}} + \mathcal{L}_{\text{proposal}}.$
\begin{equation}
\mathcal{L}_\text{all} = \mathcal{L}_{\text{RL}} + 0.5 \cdot \mathcal{L}_{\text{aux}} + \mathcal{L}_{\text{proposal}}
\end{equation}